% !TEX root =./ECE444_Practice_main.tex
%
% L14 practice set (Measurement Lab 2: Radiation Patterns). Every question is
% doable without hardware: range-geometry verification, extraction of HPBW /
% sidelobe level / front-to-back from a printed table of measured data, gain by
% comparison, cross-pol discrimination, and what a noise floor does to all of it.

\fullwidth{\textbf{Learning Objective 2.7: } I can measure the radiation pattern of an antenna and extract gain, beamwidth, sidelobe level, and polarization from the data.
\\[4pt]\normalfont\small Show your work. A \textbf{Documentation} line at the end
is required (write \textbf{None} if you did not collaborate).}

\begin{questions}

%%%%%%%%%%%%%%%%%%%%%%%%%%%%%%%%%%%%%%%%%%%%%%%%%%%%%%%%%%%%%%%%%%%%%%%%%%%%
\question \textbf{Is the range honest?} You are setting up a pattern range at
$f = 2.45\ \ghz$. The antenna under test (AUT) is a pyramidal horn with a
$24 \times 17\ \text{cm}$ aperture. The lab bench gives you a separation of
$r = 3.0\ \m$.

\begin{parts}

	\begin{minipage}[t]{0.58\linewidth}
		\part Find the wavelength and the AUT's largest dimension $D$, then
		evaluate the far-field criterion $r \ge 2D^2/\lambda$.
		\begin{solution}[1.3in]
		The largest dimension of a rectangular aperture is its diagonal, not a side.
		\eq{
			\lambda &= \frac{3\ee{8}}{2.45\ee{9}} = 0.1224\ \m = 12.2\ \text{cm} \\
			D &= \sqrt{(0.24)^2 + (0.17)^2} = 0.294\ \m \\
			\frac{2D^2}{\lambda} &= \frac{2(0.294)^2}{0.1224} = 1.41\ \m
		}
		\end{solution}
	\end{minipage}\hfill%
	\begin{minipage}[t]{0.37\linewidth}
		\raggedleft \vspace{12pt}
		$2D^2/\lambda$ = \ansbox[1.3in]{$1.41\ \m$}
	\end{minipage}

	\begin{minipage}[t]{0.58\linewidth}
		\part The far field also requires $r \ge 5D$ and $r \ge 10\lambda$.
		Evaluate both, state which criterion binds, and say whether $3.0\ \m$
		is acceptable.
		\begin{solution}[1.3in]
		\eq{
			5D &= 5(0.294) = 1.47\ \m \\
			10\lambda &= 10(0.1224) = 1.22\ \m
		}
		The binding criterion is $5D = 1.47\ \m$ --- for a physically small
		antenna the famous $2D^2/\lambda$ rule is often \emph{not} the one that
		decides. A $3.0\ \m$ range clears the binding value by a factor of
		$2.0$, so it is comfortably acceptable.
		\end{solution}
	\end{minipage}\hfill%
	\begin{minipage}[t]{0.37\linewidth}
		\raggedleft \vspace{12pt}
		binding $r$ = \ansbox[1.3in]{$1.47\ \m$}
	\end{minipage}

	\part In one or two sentences: if you ran this measurement at
	$r = 0.5\ \m$ instead, what specifically would go wrong with the recorded
	pattern, and would the measured HPBW read too wide or too narrow?

		\begin{solution}[1.1in]
		At $0.5\ \m$ the spherical wave from the source is not flat across the
		AUT aperture, so the edges of the aperture are excited with the wrong
		relative phase. That phase error broadens the main beam and fills in the
		nulls, and the sidelobes smear into the beam skirt. The measured HPBW
		reads \textbf{too wide}, and the sidelobe structure you report would not
		be the antenna's far-field pattern at all.
		\end{solution}

	\begin{minipage}[t]{0.58\linewidth}
		\part Your calibrated reference horn has a $34 \times 25\ \text{cm}$
		aperture. Recompute the far-field distance for \emph{it}, and state what
		that means for the range you just declared acceptable.
		\begin{solution}[1.3in]
		\eq{
			D &= \sqrt{(0.34)^2 + (0.25)^2} = 0.422\ \m \\
			\frac{2D^2}{\lambda} &= \frac{2(0.422)^2}{0.1224} = 2.91\ \m
			\qquad (5D = 2.11\ \m)
		}
		The \emph{reference} sizes the range, not the AUT, and $3.0\ \m$ clears
		$2.91\ \m$ by only 3\%. The gain comparison is therefore the one
		measurement standing at the edge of the far field; say so in the report,
		or lengthen the range to roughly $5.8\ \m$ for a factor-of-two margin.
		\end{solution}
	\end{minipage}\hfill%
	\begin{minipage}[t]{0.37\linewidth}
		\raggedleft \vspace{12pt}
		$2D^2/\lambda$ = \ansbox[1.3in]{$2.91\ \m$}
	\end{minipage}

\end{parts}

\newpage
%%%%%%%%%%%%%%%%%%%%%%%%%%%%%%%%%%%%%%%%%%%%%%%%%%%%%%%%%%%%%%%%%%%%%%%%%%%%
\question \textbf{Reducing a measured cut.} The table below is one measured
E-plane cut of an aperture antenna, in dBm at the receiver. The cut was
symmetric about $0\mydeg$ to within $0.3\db$, so only $\theta \ge 0$ is
tabulated. With the source switched off and every other setting unchanged, the
receiver read $-78.0\dbm$.

\begin{center}
\small
\begin{tabular}{r r @{\hspace{2em}} r r @{\hspace{2em}} r r}
\hline
$\theta$ & level & $\theta$ & level & $\theta$ & level \\
\hline
$0\mydeg$  & $-35.6$ & $18\mydeg$ & $-49.5$ & $48\mydeg$  & $-57.5$ \\
$2\mydeg$  & $-36.0$ & $20\mydeg$ & $-48.7$ & $60\mydeg$  & $-59.8$ \\
$4\mydeg$  & $-36.7$ & $22\mydeg$ & $-50.0$ & $75\mydeg$  & $-62.5$ \\
$6\mydeg$  & $-38.3$ & $24\mydeg$ & $-52.7$ & $90\mydeg$  & $-64.0$ \\
$8\mydeg$  & $-41.4$ & $28\mydeg$ & $-77.6$ & $120\mydeg$ & $-66.5$ \\
$10\mydeg$ & $-45.4$ & $32\mydeg$ & $-55.6$ & $150\mydeg$ & $-63.0$ \\
$12\mydeg$ & $-53.5$ & $36\mydeg$ & $-53.3$ & $180\mydeg$ & $-60.6$ \\
$14\mydeg$ & $-66.4$ & $40\mydeg$ & $-56.7$ &             &         \\
$16\mydeg$ & $-52.6$ & $44\mydeg$ & $-70.6$ &             &         \\
\hline
\end{tabular}
\end{center}

\begin{parts}

	\begin{minipage}[t]{0.58\linewidth}
		\part Extract the half-power beamwidth. Interpolate between samples ---
		do not snap to the nearest grid point.
		\begin{solution}[1.3in]
		The peak is $-35.6\dbm$, so the half-power level is $-38.6\dbm$. That
		falls between $6\mydeg$ ($-38.3\dbm$) and $8\mydeg$ ($-41.4\dbm$):
		\eq{
			\theta_{-3} &= 6 + 2\lp\frac{38.6-38.3}{41.4-38.3}\rp
			 = 6 + 2(0.097) = 6.19\mydeg \\
			\theta_{HP} &= 2(6.19) = 12.4\mydeg
		}
		\end{solution}
	\end{minipage}\hfill%
	\begin{minipage}[t]{0.37\linewidth}
		\raggedleft \vspace{12pt}
		$\theta_{HP}$ = \ansbox[1.3in]{$12.4\mydeg$}
	\end{minipage}

	\begin{minipage}[t]{0.58\linewidth}
		\part Find the first sidelobe level relative to the peak, and compare it
		with the canonical value for a uniformly illuminated aperture.
		\begin{solution}[1.3in]
		The first null is near $14\mydeg$; the next lobe peaks at $20\mydeg$ with
		$-48.7\dbm$:
		\eq{
			\text{SLL} = -48.7 - (-35.6) = -13.1\db
		}
		A uniformly illuminated aperture gives $-13.3\db$, so this aperture is
		essentially uniformly illuminated --- no taper. The second lobe at
		$36\mydeg$ gives $-53.3 + 35.6 = -17.7\db$, against the uniform-aperture
		value of $-17.9\db$: the same story.
		\end{solution}
	\end{minipage}\hfill%
	\begin{minipage}[t]{0.37\linewidth}
		\raggedleft \vspace{12pt}
		SLL = \ansbox[1.3in]{$-13.1\db$}
	\end{minipage}

	\begin{minipage}[t]{0.58\linewidth}
		\part Compute the front-to-back ratio.
		\begin{solution}[1.0in]
		\eq{
			\text{F/B} = -35.6 - (-60.6) = 25.0\db
		}
		\end{solution}
	\end{minipage}\hfill%
	\begin{minipage}[t]{0.37\linewidth}
		\raggedleft \vspace{12pt}
		F/B = \ansbox[1.3in]{$25.0\db$}
	\end{minipage}

	\part The operator used a $2\mydeg$ step inside the main beam. Using the
	rule of thumb for angle step, was that fine, wasteful, or too coarse?
	Answer in one sentence, with the number that justifies it.

		\begin{solution}[0.9in]
		The rule is step $\le \theta_{HP}/5 = 12.4/5 = 2.5\mydeg$, so a
		$2\mydeg$ step satisfies it with a little margin --- fine, not wasteful.
		(The $4\mydeg$ and coarser steps used past $24\mydeg$ do \emph{not}
		resolve the sidelobe peaks well, which is why the extracted sidelobe
		levels deserve a wider error bar than the beamwidth does.)
		\end{solution}

\end{parts}

\newpage
%%%%%%%%%%%%%%%%%%%%%%%%%%%%%%%%%%%%%%%%%%%%%%%%%%%%%%%%%%%%%%%%%%%%%%%%%%%%
\question \textbf{Gain by comparison, and polarization.} On the range of
Question 1 you record a peak of $-35.6\dbm$ with the AUT installed. You then
swap in the calibrated reference horn ($G_{ref} = 15.0\dbi$), change nothing
else, and record $-32.4\dbm$.

\begin{parts}

	\begin{minipage}[t]{0.58\linewidth}
		\part Find the gain of the AUT in dBi.
		\begin{solution}[1.2in]
		In dB the comparison method is a single subtraction, because transmit
		power, path loss, cable loss and source gain are common to both
		measurements and cancel:
		\eq{
			G_{AUT} &= G_{ref} + \lp P_{AUT} - P_{ref}\rp \\
			 &= 15.0 + \lp -35.6 + 32.4 \rp \\
			 &= 15.0 - 3.2 = 11.8\dbi
		}
		\end{solution}
	\end{minipage}\hfill%
	\begin{minipage}[t]{0.37\linewidth}
		\raggedleft \vspace{12pt}
		$G_{AUT}$ = \ansbox[1.3in]{$11.8\dbi$}
	\end{minipage}

	\part The procedure insists that you change \emph{nothing} between the two
	measurements. Name two things that would silently corrupt the result if you
	changed them, and say what the resulting gain error would look like.

		\begin{solution}[1.1in]
		Any change that is not common to both measurements survives the
		subtraction and becomes gain error. Examples: moving or reseating a
		cable (a fraction of a dB of extra loss appears as a fixed gain offset);
		re-peaking sloppily so the reference is measured off boresight (the
		reference reads low, so the AUT gain reads \emph{high}); changing the
		transmit level or the resolution bandwidth between the two runs; moving
		either mast so the range length differs. All of these are biases, not
		noise --- repeating the sweep will not reveal them.
		\end{solution}

	\begin{minipage}[t]{0.58\linewidth}
		\part You rotate the source antenna $90\mydeg$ about the range axis and
		re-peak: the boresight level drops to $-59.9\dbm$. Compute the
		cross-polarization discrimination.
		\begin{solution}[1.1in]
		\eq{
			\text{XPD} = P_{co} - P_{cross} = -35.6 - (-59.9) = 24.3\db
		}
		That is squarely in the healthy $20$--$30\db$ band for a linearly
		polarized antenna.
		\end{solution}
	\end{minipage}\hfill%
	\begin{minipage}[t]{0.37\linewidth}
		\raggedleft \vspace{12pt}
		XPD = \ansbox[1.3in]{$24.3\db$}
	\end{minipage}

	\begin{minipage}[t]{0.58\linewidth}
		\part Suppose instead the cross-pol peak had landed only $6.0\db$ above
		your noise floor. By how much would the measured cross-pol level
		overstate the true one, and what does that do to your reported XPD?
		\begin{solution}[1.3in]
		The reading is signal plus noise. If measured is $6.0\db$ above the
		floor, then measured$/$noise $= 10^{0.6} = 3.98$, so
		signal$/$noise $= 2.98$ and
		\eq{
			\text{error} = 10\ \mylog{\frac{3.98}{2.98}} = 1.25\db
		}
		The cross-pol level reads $1.25\db$ high, so the true XPD is
		$1.25\db$ \emph{better} than reported. XPD measured near the floor is
		therefore a \textbf{lower bound}: report it as ``XPD $\ge$'' with the
		floor quoted beside it.
		\end{solution}
	\end{minipage}\hfill%
	\begin{minipage}[t]{0.37\linewidth}
		\raggedleft \vspace{12pt}
		error = \ansbox[1.3in]{$1.25\db$}
	\end{minipage}

\end{parts}

\newpage
%%%%%%%%%%%%%%%%%%%%%%%%%%%%%%%%%%%%%%%%%%%%%%%%%%%%%%%%%%%%%%%%%%%%%%%%%%%%
\question \textbf{What the noise floor is doing to your numbers.} A measured
pattern reaches its peak at $-35.6\dbm$; with the source off, the receiver reads
$-78.0\dbm$. Levels below are quoted relative to the peak.

\begin{parts}

	\begin{minipage}[t]{0.58\linewidth}
		\part A null in a different data set measures $-25.0\db$ on a system
		whose floor sits at $-30.0\db$ relative to the peak. Estimate the true
		depth of that null.
		\begin{solution}[1.3in]
		Signal and noise add in \emph{power}, so subtract them in power:
		\eq{
			p_{true} &= 10^{-2.50} - 10^{-3.00} \\
			 &= 3.162\ee{-3} - 1.000\ee{-3} = 2.162\ee{-3} \\
			\text{null} &= 10\ \mylog{2.162\ee{-3}} = -26.7\db
		}
		The null is $1.7\db$ deeper than it measured.
		\end{solution}
	\end{minipage}\hfill%
	\begin{minipage}[t]{0.37\linewidth}
		\raggedleft \vspace{12pt}
		true null = \ansbox[1.3in]{$-26.7\db$}
	\end{minipage}

	\begin{minipage}[t]{0.58\linewidth}
		\part Now do the same subtraction for a reading of $-42.0\db$ on this
		system, whose floor is at $-78.0 + 35.6 = -42.4\db$ relative. Comment on
		the answer you get.
		\begin{solution}[1.4in]
		\eq{
			p_{true} &= 10^{-4.20} - 10^{-4.24} \\
			 &= 6.310\ee{-5} - 5.754\ee{-5} = 5.55\ee{-6} \\
			\text{null} &= 10\ \mylog{5.55\ee{-6}} = -52.6\db
		}
		The number is worthless. It is the difference of two nearly equal
		quantities: a $0.3\db$ error in either reading moves the answer by tens
		of dB. The only defensible statement is ``the null is below
		$-42\db$, floor-limited.''
		\end{solution}
	\end{minipage}\hfill%
	\begin{minipage}[t]{0.37\linewidth}
		\raggedleft \vspace{12pt}
		true null = \ansbox[1.3in]{$\le -42\db$}
	\end{minipage}

	\part You average 16 sweeps instead of 1. In one or two sentences, say what
	improves and what does not, and why.

		\begin{solution}[1.1in]
		Averaging power measurements incoherently reduces the \emph{variance} of
		the noise --- the fuzz on the trace shrinks by about $\sqrt{16} = 4$ ---
		but the \emph{mean} noise power is unchanged, so the floor does not
		move. You get a prettier trace with exactly the same dynamic range.
		Lowering the floor requires more transmit power, a narrower resolution
		bandwidth, or a quieter receiver.
		\end{solution}

	\begin{minipage}[t]{0.58\linewidth}
		\part State this system's dynamic range, then rank HPBW ($12.4\mydeg$),
		first sidelobe ($-13.1\db$) and null depth ($-42.0\db$) by how much you
		trust them, with the margin above the floor for each.
		\begin{solution}[1.4in]
		\eq{
			\text{DR} = -35.6 - (-78.0) = 42.4\db
		}
		HPBW is read $3\db$ down, roughly $39\db$ above the floor --- trust it
		first. The first sidelobe sits $29\db$ above the floor, so the noise
		adds about $0.005\db$ to it --- trust it too. The null is \emph{at} the
		floor, margin $\approx 0\db$ --- report it as a bound, never as a value.
		\end{solution}
	\end{minipage}\hfill%
	\begin{minipage}[t]{0.37\linewidth}
		\raggedleft \vspace{12pt}
		DR = \ansbox[1.3in]{$42.4\db$}
	\end{minipage}

\end{parts}

\newpage
%%%%%%%%%%%%%%%%%%%%%%%%%%%%%%%%%%%%%%%%%%%%%%%%%%%%%%%%%%%%%%%%%%%%%%%%%%%%
\question \textbf{The uncertainty paragraph you have to write.} Your report
claims $G_{AUT} = 11.8\dbi$ for the horn of Question 3, against an
aperture-formula prediction of $12.3\dbi$ (using $\eta_{ap} = 0.5$).

\begin{parts}

	\begin{minipage}[t]{0.58\linewidth}
		\part Your peak-find left the AUT $1.5\mydeg$ off boresight on a
		$12.4\mydeg$ beam. Using a parabolic approximation to the main beam near
		its peak, how much does the recorded peak read low?
		\begin{solution}[1.3in]
		Near the peak the pattern in dB is parabolic and passes through
		$-3\db$ at half the beamwidth, $6.2\mydeg$:
		\eq{
			\Delta &= -3\lp\frac{1.5}{6.2}\rp^2 = -3(0.0585) \\
			 &= -0.18\db
		}
		Small --- but it is a \textbf{bias}, always in the same direction, so
		repeating the sweep will never remove it.
		\end{solution}
	\end{minipage}\hfill%
	\begin{minipage}[t]{0.37\linewidth}
		\raggedleft \vspace{12pt}
		$\Delta$ = \ansbox[1.3in]{$-0.18\db$}
	\end{minipage}

	\begin{minipage}[t]{0.58\linewidth}
		\part A floor reflection arrives $20\db$ below the direct path. Find the
		peak-to-peak ripple it puts on your pattern, then find how far down the
		reflection would have to be to hold the ripple inside $\pm 0.2\db$.
		\begin{solution}[1.5in]
		Voltage ratio $10^{-20/20} = 0.10$; the two paths add and subtract:
		\eq{
			\text{max} &= 20\ \mylog{1.10} = +0.83\db \\
			\text{min} &= 20\ \mylog{0.90} = -0.92\db
		}
		so about $1.7\db$ peak to peak. For $\pm 0.2\db$ we need
		$20\ \mylog{1+x} = 0.2$, giving $x = 0.023$, i.e. the reflection must be
		$20\ \mylog{0.023} = -32.6\db$ below the direct path. That is what the
		absorber on the bounce point is buying you.
		\end{solution}
	\end{minipage}\hfill%
	\begin{minipage}[t]{0.37\linewidth}
		\raggedleft \vspace{12pt}
		ripple = \ansbox[1.3in]{$1.7\db$ p-p}
	\end{minipage}

	\part Combine the error sources --- reference tolerance $0.5\db$, ripple
	$0.9\db$, alignment $0.2\db$, sweep-to-sweep repeatability $0.3\db$ --- into
	one uncertainty, then write the two-sentence verdict on the $0.5\db$ gap
	between your measurement and the prediction.

		\begin{solution}[1.5in]
		Independent contributions add in quadrature:
		\eq{
			u = \sqrt{0.5^2 + 0.9^2 + 0.2^2 + 0.3^2} = \sqrt{1.19} = 1.1\db
		}
		Verdict: the measured $11.8\dbi$ and the predicted $12.3\dbi$ differ by
		$0.5\db$, well inside the $\pm 1.1\db$ uncertainty, so the measurement
		\textbf{agrees} with prediction and no physical explanation is required.
		Note also that the prediction itself is soft --- it assumed
		$\eta_{ap} = 0.5$, and real horns run anywhere from $0.4$ to
		$0.6$, which is a further $\pm 1\db$ on the number being compared
		against.
		\end{solution}

\end{parts}

\vspace{1.5em}
\noindent\textbf{Documentation:} \rule[-2pt]{0.6\linewidth}{0.4pt}

\end{questions}
